% preamble

\input{preamble}

% document

\begin{document}

\section*{$E$ is Irrational \begin{small} \begin{color}{gray} \\
(May 13, 2026) \end{color} \end{small}}

If the number $e$ is defined to be the sum of the series $\ds \sum_{n=0}^{\infty} \frac{1}{n!}$, then the error in approximating $e$ by the partial sum $\ds 1+1+1/2+\cdots+1/n!$ is 
\begin{align*}
 0 &< e - \lp 1 + \frac{1}{1!} + \cdots + \frac{1}{n!} \rp \\
 &= \frac{1}{(n+1)!} + \frac{1}{(n+2)!} + \cdots \\
 &= \frac{1}{(n+1)!} \lp 1 + \frac{1}{n+2} + \frac{1}{(n+2)(n+3)} + \cdots \rp \\
&\leq \frac{1}{(n+1)!}  \sum_{k=0}^{\infty} \lp \frac{1}{n+2} \rp^k \\
&= \frac{n+2}{(n+1)!(n+1)}
\end{align*}
The series converges quickly. Also, if $\ds e = a/b$ for positive integers $a$ and $b$, then the number 
\[
  x = b! \lp e - \lp 1 + \frac{1}{1!} + \cdots + \frac{1}{b!} \rp\rp
\] 
is an integer, and
\[
  0< x \leq \frac{b!(b+2)}{(b+1)!(b+1)} = \frac{b+2}{(b+1)^2} <1
\]
which is a contradiction.

\end{document}
